\section*{DETECTOR GEIGER-MULLER}

Dentro de los instrumentos de monitoreo se encuentran el monitor con Detector Geiger y el monitor con Detector de Centelleo.

El detector Geiger-Müller (GM) mide directamente el nivel de radiación al convertir la radiación ionizante en pulsos de corriente eléctrica y generalmente tiene tres escalas:
\begin{itemize}
    \item 0-10, 0-100, 0-1000 mR/hr.
\end{itemize}

Son aparatos confiables y en el caso de la radiografía industrial son indispensables para un trabajo seguro.

El botón de encendido tiene cinco posiciones: Apagado, baterías (para chequear carga), x10, x100, x1. La lectura de la aguja en la escala graduada de 0-10 se multiplica por 1, 10 o 100 según en donde se encuentre el selector.

Por ejemplo si la aguja estuviera marcando 2 mR/hr y el selector estuviera en la posición x10 significaría que la lectura correcta es de 20 mR/hr.

\subsection*{Ventajas y desventajas de estos instrumentos}
\subsubsection*{Ventajas}
\begin{itemize}
    \item Miden el nivel de radiación al cual se está exponiendo.
\end{itemize}

\subsubsection*{Desventajas}
\begin{itemize}
    \item No registra las dosis que ha recibido el individuo.
\end{itemize}

\section*{EL MONITOR CON DETECTOR DE CENTELLEO}

Se basa en la detección por medio de los destellos luminosos producidos por la radiación nuclear en ciertos materiales.

La ventaja de este tipo de instrumentos es que posee:
\begin{itemize}
    \item Mayor eficiencia de conteo.
\end{itemize}

\chapter{FILOSOFÍA DE LA PROTECCIÓN RADIOLÓGICA}

Tanto en lo que se refiere a trabajadores ocupacionalmente expuestos como para el público se deberá procurar que las dosis se mantengan en un nivel mínimo tomando en cuenta consideraciones de riesgo-beneficio.

Esto es contemplado por el criterio ``ALARA'' (As Low As is Reasonable Achievable) que significa "Tan bajo como razonablemente pueda alcanzarse".

Todas las actividades que involucren exposición a la radiación deberán estar gobernadas por este criterio y la evaluación ALARA de una cierta actividad podrá servir como base en la toma de decisiones.
